% Generated draft: requires mathematical and source review.
\documentclass[11pt]{article}
\usepackage[margin=25mm]{geometry}
\usepackage{fontspec}
\setmainfont{Latin Modern Roman}
\usepackage{amsmath,amssymb,mathtools,unicode-math}
\setmathfont{Latin Modern Math}
\usepackage{enumitem,needspace,xurl,hyperref}
\setlength{\parindent}{0pt}
\setlength{\parskip}{5pt}
\setlength{\emergencystretch}{3em}
\newcommand{\N}{\mathbb{N}}
\newcommand{\Z}{\mathbb{Z}}
\newcommand{\Q}{\mathbb{Q}}
\newcommand{\R}{\mathbb{R}}
\newcommand{\C}{\mathbb{C}}
\newcommand{\F}{\mathbb{F}}
\newcommand{\A}{\mathbb{A}}
\newcommand{\PP}{\mathbb{P}}
\newcommand{\E}{\mathbb{E}}
\newcommand{\eps}{\varepsilon}
\newcommand{\dd}{\,\mathrm{d}}
\begin{document}
\section*{UnsolvedMath ID: 701312}
\textbf{Source problem code:} catalogue entry
\subsection{Definitions and mathematical statement}
Let \(G\) be an arbitrary group, with identity \(1\). Use the commutator convention \([x,y]=x^{-1}y^{-1}xy\), and define iterated commutators recursively by
\[
[x,{}_1 y]=[x,y],\qquad [x,{}_{n+1}y]=[[x,{}_n y],y]\quad(n\geq 1).
\]
An element \(a\in G\) is a \emph{right \(4\)-Engel element} if
\[
\forall g\in G,\qquad [a,{}_4 g]=1.
\]
A group is \emph{locally nilpotent} if each of its finitely generated subgroups is nilpotent. The \emph{Hirsch--Plotkin radical} \(\operatorname{HP}(G)\) is the largest normal locally nilpotent subgroup of \(G\).

The question is whether
\[
\exists G\ \exists a\in G\quad
\left(\bigl(\forall g\in G,\ [a,{}_4 g]=1\bigr)
\ \land\ a\notin\operatorname{HP}(G)\right).
\]
No finiteness assumption or other restriction on \(G\) is imposed in the supplied statement.
\subsection{Short English statement}
Does there exist a group with a right \(4\)-Engel element lying outside its Hirsch--Plotkin radical, that is, outside its largest normal locally nilpotent subgroup?
\subsection{Sources}
\begin{itemize}
\item Local UnsolvedMath registry, supplied record with permanent ID \(701312\), exact source code \texttt{catalogue entry}, titled \emph{Kourovka Notebook Problem 17.10}. The registry labels the problem open; this status is supplied metadata and has not been independently verified.
\item \emph{The Kourovka Notebook (Unsolved Problems in Group Theory)}, Problem 17.10, proposed by A. Abdollahi in 2010. The supplied background identifies the 17th issue (2010) and the 21st issue edition (2026), \texttt{arXiv:1401.0300v48}. The supplied collection description names E. I. Khukhro and V. D. Mazurov as editors.
\end{itemize}
\end{document}
